\documentclass{article}
\usepackage[utf8]{inputenc}
\usepackage{graphicx}
\usepackage{listings}
\usepackage{color}
\usepackage{geometry}
\geometry{a4paper, margin=1in}

\definecolor{codegray}{rgb}{0.5,0.5,0.5}
\definecolor{codepurple}{rgb}{0.58,0,0.82}
\definecolor{backcolour}{rgb}{0.95,0.95,0.92}

\lstdefinestyle{mystyle}{
    backgroundcolor=\color{backcolour},   
    commentstyle=\color{codegray},
    keywordstyle=\color{blue},
    numberstyle=\tiny\color{codegray},
    stringstyle=\color{codepurple},
    basicstyle=\ttfamily\footnotesize,
    breakatwhitespace=false,         
    breaklines=true,                 
    captionpos=b,                    
    keepspaces=true,                 
    numbers=none,                    
    showspaces=false,                
    showstringspaces=false,
    showtabs=false,                  
    tabsize=2
}
\lstset{style=mystyle}

\title{Assignment 1 Report: Simple OpenGL Project}
\author{Group Report}
\date{\today}

\begin{document}

\maketitle

\section{Introduction}
This report details the implementation of a simple OpenGL application. The project successfully loads a 3D model (Trophy Table) downloaded from a 3D asset provider, and renders it using custom vertex and fragment shaders. The application features user controls to modify matrices, lighting conditions, and materials.

\section{Implementation Details}

\subsection{Vertex and Fragment Shaders (Phong Model)}
The application implements the standard Phong reflection model directly within the \texttt{fragment.glsl} file. The shader calculates three main components of lighting: Ambient, Diffuse, and Specular.

\begin{itemize}
    \item \textbf{Ambient Lighting:} Simulates indirect light bounding off surfaces so that no shadow is completely black.
    \item \textbf{Diffuse Lighting:} Simulates the directional impact of the light object, making parts facing the light brighter.
    \item \textbf{Specular (Phong) Lighting:} Simulates the bright spot/reflection of a light on objects.
\end{itemize}

Below is the snippet from the fragment shader calculating the Phong reflections:

\begin{lstlisting}[language=C++]
vec3 norm = normalize(Normal);
vec3 viewDir = normalize(viewPos - FragPos);

// Diffuse
float diff = max(dot(norm, lightDir), 0.0);
diffuse += intensity * diff * color * material.diffuse * texColor;

// Specular (Phong)
vec3 reflectDir = reflect(-lightDir, norm);  
float spec = pow(max(dot(viewDir, reflectDir), 0.0), material.shininess);
specular += intensity * spec * color * specMask;
\end{lstlisting}

Different parts of the mesh are given different material characteristics (like glass, metal, and wood) by passing different \texttt{material.shininess} and \texttt{material.alpha} values via uniforms. The code also implements logic to discard pixels for transparent logic separating opaque passes and blending passes.

\subsection{Loading a 3D Model (VAO \& VBO)}
The project loads complex \texttt{.obj} files using the Assimp (Asset Importer) Library. In the codebase, the \texttt{Model} and \texttt{Mesh} classes handle the abstraction. Once vertex data (Position, Normal, Texture Coordinates) is parsed, it is securely bound to the GPU leveraging Vertex Array Objects (VAO) and Vertex Buffer Objects (VBO).

\begin{lstlisting}[language=C++]
glGenVertexArrays(1, &VAO);
glGenBuffers(1, &VBO);
glGenBuffers(1, &EBO);

glBindVertexArray(VAO);
glBindBuffer(GL_ARRAY_BUFFER, VBO);
glBufferData(GL_ARRAY_BUFFER, vertices.size() * sizeof(Vertex), 
             &vertices[0], GL_STATIC_DRAW);
\end{lstlisting}

This approach ensures the pipeline only copies data to the graphics card once, which is highly efficient. In the rendering loop, \texttt{glBindVertexArray(VAO)} is called before invoking \texttt{glDrawElements}.

\subsection{Model, View, and Projection Matrices Setup}
The 3D space transformations are thoroughly demonstrated with interactive key-bindings:

\begin{itemize}
    \item \textbf{Model Matrix:} Controls the position, scale, and rotation of the object. By pressing the \textbf{'R'} key, an auto-rotation feature modifies the model matrix over time frame-by-frame using \texttt{glm::rotate}.
    \item \textbf{View Matrix:} Simulates a moving camera. Users can press \textbf{'V'} to cycle between three defined camera offsets (Cinematic, Side, and Top perspectives) achieved using \texttt{glm::lookAt()}.
    \item \textbf{Projection Matrix:} Simulates the lens of the camera. Pressing \textbf{'P'} toggles the application between a perspective projection (simulating human eyes via \texttt{glm::perspective()}) and an orthographic projection (ignoring depth variance via \texttt{glm::ortho()}).
\end{itemize}

\section{Screenshots}
\textit{(Please insert the final screenshots of the running app here. Included pictures should demonstrate the Orthographic vs Perspective projections, Textures enabled/disabled, and different view angles.)}

\end{document}